% =====================================================================
%  sec05 -- The [T] lemma table
% =====================================================================
\section[The eight lemmas, one sentence each]{The eight \LT{} lemmas, one sentence each}
\label{sec:lemmas}

Everything below is proved, and every proof names its file. The list is short on purpose: after
removing the results that fail the free-field acid test (\S\ref{sec:dead}, item~9) and the results
that were withdrawn (\S\ref{sec:errata}), this is what the coarse-mode wave actually established.
The lemmas are listed because they are true, not because they support the route; two of them
(lemma~4, lemma~6) are in fact the reason two of the dead ends are dead.

{\footnotesize
\begin{longtable}{@{}L{0.5cm}L{5.0cm}L{4.2cm}L{4.2cm}@{}}
\toprule
\textbf{\#} & \textbf{Statement} & \textbf{One line of content} & \textbf{Source} \\
\midrule
\endhead
1 & $\Phi''\succeq(N/12)\id$ in every direction; pointwise
    $\Phi''\ge N/12+a_0^{2}/120$
  & the floor is isotropic, sharp (attained only at $a_0=0$), and \emph{grows} with the coarse
    amplitude; no $a_0$-weakening
  & \code{P16-PHIQUANT.md} \S0; \code{ch05-theorems.tex}:207--209 \\
2 & $E[H_{vv}]\ge0$ for $\theta\le\pi$
  & $H_{vv}=\beta Q$ with $Q=\tfrac12\cos(\theta/2)g(u)\ge0$ \emph{pointwise}, so the curvature
    side is free in this window
  & \code{P05-CINF.md} \S0 item 1 \\
3 & $E[H_{vv}]=E[G_v^{2}]$ (Stein)
  & the mean Hessian equals the covariance of the gradient: half of an identity, not a spectral
    condition
  & \code{Q15-SPECTRAL.md} \S2 \\
4 & $E[\Phi'']=E[\Phi'^{2}]+E[G\Phi']$ (Stein/IBP)
  & exact, with no boundary term in the log chart; equivalently
    $\Var(\Phi')=E[\Phi'']-(E\Phi')^{2}-\Cov(G,\Phi')$
  & \code{S04-VARPHI.md} verdict, first display \\
5 & $H=\tfrac12\cot(\theta/2)G$ at $a=0$, pointwise
  & $G$ and $H$ are the same function of $u$ up to a constant; residual $10^{-16}$; at $\theta=\pi$
    it degenerates to $H\equiv0$
  & \code{Q04-WEYL.md} \S0; \code{R03-CERT-THETA.md} \S1 \\
6 & $\Var(w)\le\tfrac12\langle w\rangle$ and
    $\langle w\rangle-2\Var\ge2\langle w\rangle^{2}\ge0$
  & two nested bounds, three independent derivations, valid for all $\beta$ and all $c$; the
    measured ratio $\langle w\rangle/\Var$ has minimum $32.3177$ at $\gamma=0$
  & \code{R02-CRIT.md} \S0; \code{R05-WCF.md} \S0, \S6; \code{R09-FLUCT.md} \S4 \\
7 & skeleton identity
    $\Var(F)-E[\kappa]=-\Gpp=-m_{G'}-E[\Phi'']+\Var(\Phi')+\Cov(G,\Phi')$
  & pure algebra joining the split, the Ward form and the $\Phi$ decomposition; its numerical
    content is a check (residual $\le7\times10^{-6}$ on a six-point $a$-grid; $0.06\%$ on $L=2$)
  & \code{P21-REBUILD.md} \S0; \code{ch19-h1-status.tex}:474--491 \\
8 & $H=K_0G+\beta\Delta$ with $K_0=\tfrac12\cot(\theta/2)$ and
    $\Delta=tD-(c^{2}/4t)A=O(a)$ odd
  & the exact replacement for the false affinity $H=K(a)G+C(a)$: same operator, different argument,
    one explicit function apart
  & \code{R07-AFFINE.md} \S0 \\
\bottomrule
\end{longtable}
}

\subsection{Lemma 4 and what it forces}

Lemma 4 is the sharpest small result in the list because it is an \emph{identity}, not a bound: it
converts a variance into a mean plus a covariance, so no estimate can improve it. Its first
corollary is the sentence that kills the ``split the problem in two halves'' strategy
(\S\ref{sec:dead}, item~7): the $\Phi$ sector's contribution is exactly
\begin{equation}
E[\Phi'']-\Var(\Phi')=(E\Phi')^{2}+\Cov(G,\Phi'),
\label{eq:cor4}
\end{equation}
a square plus a cross term, and on the $a=0$ slice $E[\Phi'\,|\,0]=0$ by the $X\to-X$ symmetry, so
the $\Phi$ sector is \emph{exactly neutral} --- neither positive nor negative, with zero free
margin. \LT{} Its second corollary is the instrument that caught an error in the $\mathrm{SU}(3)$
computation: a reported $\Var(\Phi')=6.81$ exceeding a reported $E[\Phi'']=6.70$ is impossible
under \eqref{eq:cor4} with a $\Cov\ge0$, and since the pair used there has
$\Cov=-0.11$ it is still impossible by a factor of two; hence the pair
$(\Phi_{O09},|\Delta|^{2}e^{-S_W})$ is not IBP-compatible and the number is not a $\Var(\Phi')$ of
any compatible measure. The correct pairing gives $E[\Phi'^{2}]=E[\Phi'']=8.00000000$ exactly, i.e.\
$C=1.000000$. \LT{}

\src{\code{S04-VARPHI.md} verdict (all three corollaries verbatim, including
``\emph{exactly} $0$ (neutral)'', the ``$6.81$ is not any compatible measure's $\Var(\Phi')$''
sentence, and the exact $8.00000000$); \code{Q01-PHI-ONLY.md} \S0 (\S3.1: on the $\Phi$-only
measure $\Var(\Phi')\le E(\Phi'')$ holds but with residual exactly $(E\Phi')^{2}=0$);
\code{Q03-SEPARATE.md} \S0.2 (the numerical confirmation: at $a=0$, $C=2\Cov=0.0000$, and the
$\Phi$ half is the only block that holds on all $28$ grid points).}

\subsection{Lemma 3 does not give a spectral condition}

Lemma 3 is often quoted as ``$H1$ is a spectral statement''. It is not, and the difference
matters. What the lemma says is that $E[H_{vv}]=E[G_v^{2}]$ --- the mean Hessian equals the
second moment of the force --- which is one half of the covariance identity for the gradient. The
only \emph{named} spectral equivalence in the neighbourhood is
$\lambda_{\max}\le1\iff\Lambda\preceq\id$, i.e.\ the Brascamp--Lieb variance inequality /
Bakry--\'Emery $\Gamma_2\ge0$ Gaussian comparison; on the real (non-quadratic) $S_W$ it simply
fails, with $\lambda_{\max}=\beta=2.5$ and a $10.5\%$ negative-mode fraction at $\beta=2.5$. So
the spectral route is empty on the hypothesis side: the condition degenerates to $0\ge0$, while
the target inequality itself is true by lemma 3. \LN{}

\src{\code{Q15-SPECTRAL.md} \S4(b) and \S4 net conclusion (``$\Var(G)\le E[H]$ is \emph{half of an
identity} \ldots{} not a spectral condition''; $\{\lambda_k\}=\{\beta\cos\varphi\}$,
$E[\lambda]=1.268$, $\lambda_{\max}=2.5$, $P(\lambda<0)=0.105$); \code{G4-GENERAL-BOUND.md} \S2
Lemma G2 (the flat-chart counterpart, where the inequality becomes $\Wg\ge0$).}

\subsection{Lemma 8 replaces a false affinity with an exact splitting}

The corpus had been using $H=K(a)G+C(a)$ with $R^{2}>0.9992$ as if it were structure. It is not:
the Wronskians $W(G,H)=GH'-G'H$ and $W(G',H')$ are both non-zero, of first order $O(a)$
($-7.21\times10^{-4}$ and $-8.46\times10^{-6}$ at $a=0.3$, $\rho=2.5$), so the two functions are
neither proportional nor affine; and $H-KG$ has no constant term in the basis
$\{\sin(\rho/2),\cos(\rho/2),1\}$ (whose Wronskian is $-1/8$), so a strict affinity would force
$C=0$, i.e.\ strict proportionality, which the ratio computation also kills
($Q_H/Q_G=K_0+3/a-3a/\rho^{2}$ varies with $\rho$, and the two channels differ by a factor $30$:
$K_c=10.356$ against $K_s=0.308$ at $a=0.3$). \LN{} What \emph{is} exact is lemma 8, and it says
the deviation is not a constant but an explicit odd function: the high $R^{2}$ is a measure of
near-collinearity of smooth functions on a short interval, not a structural identity. \LT{}

\src{\code{R07-AFFINE.md} \S0 (the Wronskian verdict, the constant-term argument, the two-channel
factor $30$, and \code{eq:H=K0G+betaDelta}); \code{R01-HG.md} (the original fit, its $R^{2}$ values
$0.9999871/0.9998552/0.9991783$, and the closed forms $K=0.3929/0.2715/0.1768$,
$C=-0.0244/-0.0624/-0.1030$); \code{ch19-h1-status.tex}:498--506 (the same audit in the v2
register).}

\subsection{The list is what survives the acid test}

Of the four candidate quantities $\{\Cov(G,\Phi'),\Var(\Phi'),m_{G'},E[\Phi'']-N/12\}$ that the
corpus had been treating as interchangeable, the three that involve $\Phi$ alone \emph{fail} the
free-field test: $\Var(\Phi')$ tends to $2.341\times10^{-3}$, which is $65\%$ of its interacting
value rather than zero, and $E[\Phi'']-N/12$ is \emph{identical} at $\beta=2.5$ and
$\beta=10^{-6}$ because at $\theta=\pi$ the measure is Haar for every $\beta$. What survives is
$\{m_{G'},-\Cov\}$, both containing $G=\partial_vS_W$ --- the only object in the list carrying a
power of $\beta$. \LT{}/\LE{} This is the reason the eight lemmas above are stated as they are, and
it is the reason lemma 6 (which is about $w$, a $\beta$-free shape function, and its measure)
survives while the four-name picture does not: lemma 6 is a statement about the \emph{measure},
and it is sharp, but it is not a statement about dynamics either.

\src{\code{P21-REBUILD.md} \S1--\S2 (the acid test on the four candidates, the free-field values,
and the surviving pair); \code{ch19-h1-status.tex}:485--491 (verbatim); \code{P13-ACID.md}:12
(the acid-test criterion itself, and the $\beta_0$ failure discussed in \S\ref{sec:dead}).}
